\documentclass[11pt]{article}

\usepackage[margin=1in]{geometry}
\usepackage{listings}
\usepackage{xcolor}
\usepackage{hyperref}
\usepackage{amsmath}
\usepackage{parskip}

\hypersetup{
  colorlinks=true,
  linkcolor=blue!60!black,
  urlcolor=blue!60!black
}

\lstset{
  basicstyle=\ttfamily\small,
  frame=single,
  breaklines=true,
  showstringspaces=false,
  columns=fullflexible,
  keywordstyle=\color{blue!70!black},
  commentstyle=\color{green!40!black},
  stringstyle=\color{red!60!black}
}

\title{CS201 Fall 2025 --- Lecture Notes}
\author{Clarissa Littler}
\date{Fall 2025}

\begin{document}
\maketitle

\section*{Preface}

This document is a running narrative of what we covered in each class
meeting this term. Each section pairs the agenda for the day with the
little code samples we actually typed up together, so if you missed a
lecture or want to go back and re-read, the whole arc of a day should be
here in one place. The source files themselves live in the
\texttt{lecture$N$/} subdirectories of the class repo, and all of the
snippets quoted below are verbatim from those files.

The numbering reflects the lecture sequence as it actually ran ---
some numbers are skipped (those days were holidays, midterms, or
discussion-only sessions). After the numbered sequence I've included
the file-I/O day, the sockets day, and a short appendix on the
pointer-pyrotechnics examples.

\tableofcontents

\section{Lecture 1 --- Welcome, C vs.\ C++, Bit-Munging, and Two's Complement}

This was an orientation day plus the first run at the parts of C that
trip people up coming from C++, plus our opening foray into bitwise
operations and integer representation.

\subsection*{Logistics}

The class is my own personal admixture of four things:
\begin{itemize}
  \item C
  \item x86-64 assembly
  \item Concurrency
  \item Operating systems theory and practice
\end{itemize}
Our official hardware is \texttt{cslinux.pcc.edu}; if you can't log
in yet, the stopgap is \texttt{onlinegdb.com}. The repo is where the
textbook lives, where these Agenda files come from, and where every
bit of code we write together gets dropped.

\subsection*{C: what's different from C++?}

The smaller surprises:
\begin{itemize}
  \item \texttt{struct} is its own keyword and you usually have to
    write \texttt{struct Foo} when referring to one.
  \item No native \texttt{bool} --- nonzero is truthy, zero is falsy,
    and that's that.
\end{itemize}
The bigger surprises are the I/O functions: \texttt{printf} and
\texttt{scanf} instead of streams. \texttt{firstC.c} is our first
worked example, showing both \texttt{malloc} on an array and
\texttt{printf} as the way you talk to the world:

\begin{lstlisting}[language=C]
#include "stdio.h"
#include "stdbool.h"
#include "stdlib.h"

struct Rect {
  int width;
  int height;
};

int main(){
  int* arr1 = malloc(100*sizeof(int));

  for(int i=0; i<100; i++){
    *(arr1 + i) = i*i;
    printf("The %d th element is %d\n",i,arr1[i]);
  }

  free(arr1);
  return 0;
}
\end{lstlisting}

Two things to soak in here. First, \texttt{*(arr1 + i)} and
\texttt{arr1[i]} are literally the same thing --- indexing is just
pointer arithmetic in disguise. We'll lean on that all term. Second,
\texttt{malloc} returns a pointer and needs a size in bytes; unlike
\texttt{new T} in C++, it doesn't know your type.

(Side note that will come back to bite us in the audit: I keep
writing \texttt{\#include "stdio.h"} with quotes instead of
\texttt{<stdio.h>}. It works on our setup because the system headers
end up resolved either way, but the angle-bracket form is the
correct convention for system headers --- I'll start using it more
strictly later.)

\subsection*{Bitwise operations}

We went through the usual cast: \texttt{\&} (AND), \texttt{|} (OR),
\texttt{\^{}} (XOR), \texttt{\textasciitilde{}} (NOT), and the
shifts \texttt{<<} and \texttt{>>}. The classic XOR trick is
``XOR-ing a bit with 1 flips it, XOR-ing with 0 leaves it alone,''
which is the entire content of \texttt{toggleBit} in
\texttt{bitmunger1.c}:

\begin{lstlisting}[language=C]
void printBits(int num1){
  printf("%d in bits is: ",num1);
  for(int i=31; i>=0;i--){
    printf("%d",(num1>>i)&1);
  }
  printf("\n");
}

void toggleBit(int* num,int place){
  *num = (*num) ^ (1 << place);
}
\end{lstlisting}

The \texttt{printBits} function is itself a nice little exercise in
bit-twiddling: shift the bit you want down to position 0, then mask
with \texttt{\& 1} to get just that bit. The driver loops, prints
the integer in both signed and unsigned interpretations, and lets
you flip bits interactively until you type \texttt{-1}.

\subsection*{Two's complement, derived}

A plain $n$-bit unsigned integer encodes
\[
  \sum_{i=0}^{n-1} b_i \cdot 2^i,
\]
so all-ones in 32 bits is $2^{32}-1$. In two's complement, the top
bit's weight goes negative:
\[
  -b_{n-1} \cdot 2^{n-1} + \sum_{i=0}^{n-2} b_i \cdot 2^i.
\]
That's the whole magic. The largest signed 32-bit value is
$2^{31}-1$, the smallest is $-2^{31}$, and all-ones is always $-1$
because $-2^{n-1} + (2^{n-1} - 1) = -1$.

\section{Lecture 2 --- Floating Point, IEEE-754, \texttt{malloc} on Structs}

Today we changed our bit-munger to look at floats, talked through
the IEEE-754 standard, then pivoted to \texttt{malloc} on structs
and pointers in C.

\subsection*{Floats and the bit-munger}

\texttt{bitmunger2.c} is structurally the same as last time's
\texttt{bitmunger1.c}, but it reinterprets the bits of a
\texttt{float} via a pointer cast. The little spaces it prints
(between bits 31/30 and 23/22) are previewing the IEEE-754
sign / exponent / fraction split:

\begin{lstlisting}[language=C]
void printBits(int* num1){
  printf("in bits is: ");
  for(int i=31; i>=0;i--){
    printf("%d",(*num1>>i)&1);
    if(i == 31 || i == 23){
      printf(" ");
    }
  }
  printf("\n");
}

int main(){
  float num1;
  printf("Choose a starting number: ");
  scanf("%f",&num1);
  int choice;

  do {
    printf("The number %f ",num1);
    printBits((int*)&num1);
    printf("Enter a bit to flip: ");
    scanf("%d",&choice);
    toggleBit((int*)&num1,choice);
  }while (choice != -1);
}
\end{lstlisting}

The \texttt{(int*)\&num1} cast is the load-bearing trick: take the
address of the float, lie to the compiler about its type, and now
we can twiddle bits with integer operations and watch the
floating-point value change.

\subsection*{IEEE-754, formally}

A 32-bit float has:
\begin{itemize}
  \item 1 sign bit
  \item 8 exponent bits (stored with bias 127)
  \item 23 fraction bits
\end{itemize}
For normal values, the value encoded is
\[
  (-1)^s \cdot 2^{e - 127} \cdot \left(1 + \sum_{k=1}^{23} f_k \cdot 2^{-k}\right).
\]
Edge cases:
\begin{itemize}
  \item Exponent all 1s, fraction zero $\Rightarrow$ $\pm\infty$.
  \item Exponent all 1s, fraction nonzero $\Rightarrow$ \texttt{NaN}.
  \item Exponent all 0s $\Rightarrow$ subnormals: the implicit
    leading 1 is dropped and the exponent is fixed at $2^{-126}$, so
    we can gracefully approach 0.
\end{itemize}

\subsection*{Where floats lie to you}

\texttt{floatround.c} is the gotcha:

\begin{lstlisting}[language=C]
double num = 0.00001;
double acc = 0;

for(int i = 0; i < 100000; i++){
  acc = acc + num;
}

if(acc != 1){
  printf("UH OH %f is not equal to 1\n",acc);
}
if((acc < 1 + 0.000001) && acc > (1 - 0.000001)){
  printf("But it's close enough!\n");
}
\end{lstlisting}

We accumulate a quantity that, in math-land, should be exactly 1.
It isn't, because $0.00001$ has no exact binary representation, and
we're piling up the rounding error 100{,}000 times. The
``but it's close enough'' check is the practical answer: never
compare floats with \texttt{==}; always use a tolerance.

The deeper reason this is unavoidable: the rationals are countably
infinite, the reals are uncountable. There is no way to cover the
reals with finitely many bits; we're forced to scatter
``breadcrumbs'' across the number line and round everything to its
nearest breadcrumb.

\subsection*{\texttt{malloc} on structs and pointer sizes}

\texttt{malloc2.c} extends \texttt{malloc} from arrays to structs:

\begin{lstlisting}[language=C]
struct Rect {
  int width;
  int height;
};

int main(){
  struct Rect* r = malloc(sizeof(struct Rect));

  r->width = 10; // (*r).width
  r->height = 20;

  printf("%d\n", r->width * r->height);

  free(r);
  return 0;
}
\end{lstlisting}

The arrow \texttt{->} is shorthand for ``dereference, then access a
field.'' \texttt{r->width} is the same as \texttt{(*r).width} ---
the parens matter because \texttt{*} has lower precedence than
\texttt{.}.

\texttt{malloc3.c} shows that \texttt{sizeof} on a struct gives
back the total bytes the compiler reserved (which can include
padding):

\begin{lstlisting}[language=C]
struct Goofus {
  int arr[100];
};

int main(){
  struct Goofus* g = malloc(sizeof(struct Goofus));
  printf("A goofus is %ld bytes big\n",sizeof(struct Goofus));
  free(g);
  return 0;
}
\end{lstlisting}

And \texttt{pointerSize.c} drives home that all pointers are the
same size --- 8 bytes on a 64-bit machine --- regardless of what
they point to:

\begin{lstlisting}[language=C]
int num1;
float num2;
long num3;
double num4;
printf("The size of an int* is: %ld\n",sizeof(&num1));
printf("The size of a long* is: %ld\n",sizeof(&num3));
printf("The size of a float* is: %ld\n",sizeof(&num2));
printf("The size of a double* is: %ld\n",sizeof(&num4));
\end{lstlisting}

\subsection*{Stacks (the data structure)}

The agenda also has a sketch for a stack data structure exercise.
We didn't fully implement it in class, but the shape is:

\begin{lstlisting}[language=C]
struct Stack {
  // an array, plus a current length
  // pick a finite max size
};

void push(struct Stack *s, int i);
int  pop(struct Stack *s);
int  length(struct Stack s);
void printStack(struct Stack s);
\end{lstlisting}

This will reappear later when we talk about the \emph{call} stack
in assembly, which is a real instance of this same data structure
managed by the CPU.

\section{Lecture 3 --- Hex, Pointer Arithmetic, Casting, and Endianness}

Today was about getting fluent with hex output, the actual
mechanics of pointer arithmetic, casting between pointer types, and
seeing endianness with our own eyes.

\subsection*{Why hex}

\texttt{whyHex.c} is a trivial program that asks \texttt{printf} to
show the same number three ways:

\begin{lstlisting}[language=C]
int num1 = 1234;
printf("As decimal: %d\n",num1);
printf("As hex: %x\n",num1);
printf("As octal: %o\n",num1);
\end{lstlisting}

The reason hex shows up everywhere --- addresses, color codes, byte
dumps --- is that two hex digits are exactly one byte
($16 \cdot 16 = 256 = 2^8$). So \texttt{0xFF} is one byte all-ones,
which as unsigned is $255$.

\texttt{printHex.c} prints the address of every element of a
\texttt{char[10]}, which lets you see the addresses going up by
exactly 1 each step (because \texttt{sizeof(char) == 1}).

\subsection*{Pointer arithmetic, made tangible}

\texttt{pointerArith.c} shows the rule you have to internalize:
pointer arithmetic scales with the \emph{pointee} type:

\begin{lstlisting}[language=C]
int num1;
printf("The address of num1 is: %p\n",(void*)&num1);
printf("The address of &num1 + 1 is: %p \n",(void*)(&num1 + 1));
\end{lstlisting}

The two addresses printed differ by 4 (because \texttt{sizeof(int)
== 4}), not by 1.

\subsection*{Casting and ``zoom level''}

\texttt{byteMunging.c} is a great little demo: cast the same array
pointer to different types and watch \texttt{+1} mean different
things:

\begin{lstlisting}[language=C]
int arr[10];
for(int i = 0; i < 10; i++){
  arr[i] = i*i;
}

printf("If we don't cast at all arr+1 is: %p\n", (void*)(arr + 1));
printf("If we cast to long, arr+1 is: %p\n", (void*)(((long*)arr) + 1));
printf("If we cast to char, arr+1 is: %p\n", (void*)(((char*)arr) + 1));
\end{lstlisting}

As an \texttt{int*}, \texttt{+1} jumps 4 bytes; as a \texttt{long*},
\texttt{+1} jumps 8; as a \texttt{char*}, \texttt{+1} jumps 1. The
``zoom level'' of a pointer is encoded entirely in the type.

\subsection*{Endianness}

\texttt{byteOrder.c} actually peers into the bytes of an integer:

\begin{lstlisting}[language=C]
int num = 0x1a2b3c4d;

unsigned char* pty = (unsigned char*)&num;
for(int i = 0; i<4; i++){
  printf("%x ",pty[i]);
}
printf("\n");
\end{lstlisting}

On our x86-64 boxes this prints \texttt{4d 3c 2b 1a} --- the
opposite order from how we wrote the literal. That's
little-endian: the least-significant byte is at the lowest address.
The cast to \texttt{unsigned char*} is the trick: once we have a
byte-pointer, pointer arithmetic moves us one byte at a time, so
we can walk the actual storage.

There's also \texttt{byteOrder.s}, which is what \texttt{gcc -S}
emits for this C program. We didn't pick through it line-by-line
yet, but it's there to look at --- the \texttt{movzbl},
\texttt{leaq}, and \texttt{call printf@PLT} are all things we'll
unpack in detail when we get to assembly proper.

\subsection*{Two's complement, redux}

The agenda revisits two's complement with worked examples for one
byte:
\begin{itemize}
  \item Unsigned: \texttt{0111 1111} = 127, \texttt{1000 0000} = 128,
    \texttt{1111 1111} = 255.
  \item Signed (two's complement): \texttt{0111 1111} = 127,
    \texttt{1000 0000} = $-128$, \texttt{1111 1111} = $-1$.
\end{itemize}
And it shows the wraparound: $127+1$ in 8 bits is \texttt{1000 0000},
which as unsigned is 128 and as signed is $-128$. That's overflow,
and it's why $2^{31}-1+1$ becomes $-2^{31}$ on a 32-bit signed int.

The agenda also flags two real-world dates this matters for: Y2K
(rolling 2-digit years from 99 to 00) and Y2038 (the moment 32-bit
signed Unix timestamps overflow). Same arithmetic, different
calendar.

\section{Lecture 5 --- Your First Assembly: \texttt{exit}, \texttt{add}, and Looking at What \texttt{gcc -S} Emits}

(There was no Lecture 4 directory --- we shifted scheduling.)
Today was the big jump from C to assembly. We talked about what a
processor actually \emph{is}, then wrote the smallest possible
assembly programs against the bare Linux syscall interface.

\subsection*{What is a processor?}

\begin{itemize}
  \item A CPU fetches instructions from memory and does what they
    say. Most instructions are one clock cycle (with a giant
    asterisk).
  \item All modern CPUs have a \emph{clock} --- an oscillator that
    drives signals to happen synchronously.
  \item The set of instructions a CPU understands is its
    \emph{instruction set architecture} (ISA). Common ones:
    x86 (legacy 32-bit), x86-64 (modern 64-bit, what we're using),
    ARM64 (Apple Silicon, phones), RISC-V (open hardware).
  \item Asking humans to write raw byte sequences is awful. So we
    invented \emph{assembly languages} --- human-readable shorthand
    for the instruction set, defined by an \emph{assembler}. We're
    using GNU \texttt{as}, which speaks AT\&T syntax.
\end{itemize}

\subsection*{Pipeline: compiler $\to$ assembler $\to$ linker}

\begin{itemize}
  \item The \emph{compiler} turns C into assembly.
  \item The \emph{assembler} turns assembly into raw machine
    instructions in an object file.
  \item But an object file isn't yet executable --- the \emph{linker}
    pulls in library code and OS boilerplate to produce the actual
    binary.
\end{itemize}

You can stop after the compiler stage with \texttt{gcc -S file.c},
which spits out \texttt{file.s}. We'll use that all term to peek
at what C compiles to.

\subsection*{Registers}

\begin{itemize}
  \item Registers are the CPU's tiny, blazing-fast scratchpad:
    16 general-purpose registers in x86-64, each 64 bits.
  \item In AT\&T syntax all register names start with a \texttt{\%}.
  \item Some registers have conventional uses (we'll fill that table
    in next time); others (\texttt{\%r8}--\texttt{\%r15}) you can
    do whatever you want with.
\end{itemize}

\subsection*{The smallest program that doesn't crash}

\texttt{smallest.s} is the smallest assembly program that runs
without segfaulting:

\begin{lstlisting}
        .text
        .global _start

_start:
        mov $60, %rax
        mov $-1, %rdi
        syscall
\end{lstlisting}

That's it. The Linux \texttt{exit} syscall is number 60, the exit
status goes in \texttt{\%rdi}, and \texttt{syscall} actually pulls
the trigger. Run it and \texttt{echo \$?} --- you'll see 255
(which is how the shell reads $-1$ as an unsigned byte).

A few notes:
\begin{itemize}
  \item AT\&T immediates are prefixed with \texttt{\$}: \texttt{\$60}
    is the constant 60, while \texttt{60} on its own would mean
    ``the contents of memory address 60'' (which would be a
    segfault).
  \item \texttt{.text} is the section where code lives.
  \item \texttt{.global \_start} tells the linker that \texttt{\_start}
    is the entry point. There's no \texttt{main} runtime here ---
    the kernel jumps straight to \texttt{\_start}.
\end{itemize}

\subsection*{\texttt{add} as a binary operation}

\texttt{smaller.s} adds two numbers and feeds the sum back as the
exit status:

\begin{lstlisting}
        .text
        .global _start

_start:
        mov $1, %r8
        mov $2, %r9
        add %r8, %r9
        mov %r9, %rdi
        mov $60, %rax
        syscall
\end{lstlisting}

The crucial syntax point that will bite you all term:
\textbf{in AT\&T syntax, \texttt{add S, D} computes \texttt{D = D + S}.}
The destination is on the \emph{right}. So \texttt{add \%r8, \%r9}
leaves $1+2=3$ in \texttt{\%r9}.

\subsection*{Skimming what \texttt{gcc -S} emits}

For comparison, we wrote a tiny C program \texttt{adding.c} that
adds two numbers via a function:

\begin{lstlisting}[language=C]
int adder(int n1, int n2){
  int res;
  res  = n1 + n2;
  return res;
}

int main(){
  int n1 = 10;
  int n2 = 20;
  printf("%d\n",adder(n1,n2));
  return 0;
}
\end{lstlisting}

\texttt{gcc -S} produces \texttt{adding.s}, which is noisy with
\texttt{.cfi\_*} directives, frame setup, and PLT calls, but the
actual work is recognizable:

\begin{lstlisting}
adder:
        pushq   %rbp
        movq    %rsp, %rbp
        movl    %edi, -20(%rbp)
        movl    %esi, -24(%rbp)
        movl    -20(%rbp), %edx
        movl    -24(%rbp), %eax
        addl    %edx, %eax
        movl    %eax, -4(%rbp)
        movl    -4(%rbp), %eax
        popq    %rbp
        ret
\end{lstlisting}

\texttt{\%edi} and \texttt{\%esi} are the 32-bit halves of
\texttt{\%rdi} and \texttt{\%rsi} --- the first two integer
arguments. The compiler spills them onto the stack
(\texttt{-20(\%rbp)}, \texttt{-24(\%rbp)}), reloads them, adds, and
returns the result in \texttt{\%eax}. We'll formalize this calling
convention next time.

\texttt{hello.c} is just \texttt{printf("Hello world\textbackslash n");}
--- the placeholder for the C version of what we'll build in pure
assembly later.

\section{Lecture 6 --- \texttt{cmp}, \texttt{jmp}, Loops, and Globals}

With one assembly instruction (\texttt{add}) under our belts, we
filled in control flow: \texttt{cmp}, the conditional jumps, and a
real loop.

\subsection*{Quick recap}

\texttt{smallredux.s} and \texttt{arith.s} are warm-ups from last
time:

\begin{lstlisting}
_start:
        mov $5,%r8
        add $10,%r8             # r8 = r8 + 10
        add %r8,%r8             # r8 = r8 + r8
        mov %r8,%rdi

        mov $60,%rax
        syscall
\end{lstlisting}

Exit status is 30. Adding a register to itself is the cheap way to
double a value.

\subsection*{\texttt{cmp} and the conditional jumps}

\texttt{cmp S, D} computes \texttt{D - S} and sets flags but
doesn't store the result:
\begin{itemize}
  \item ZF (zero flag): was the result zero?
  \item SF (sign flag): was it negative?
  \item OF (overflow flag): did signed overflow happen?
  \item CF (carry flag): did unsigned carry happen?
\end{itemize}
The conditional jumps then read those flags. Critical AT\&T
gotcha: after \texttt{cmp \%r8, \%r9}, the jump \texttt{jl} means
\texttt{\%r9 < \%r8} (because the implicit comparison is
\texttt{\%r9 - \%r8}). Trips me up every time.

\texttt{cmp1.s} branches on equality:

\begin{lstlisting}
_start:
        mov $2,%r8
        mov $2,%r9
        cmp %r8,%r9
        je equal
        mov $0,%rdi             # 0 if not equal
        jmp finish

equal:
        mov $1,%rdi             # 1 if equal

finish:
        mov $60,%rax
        syscall
\end{lstlisting}

\texttt{cmp2.s} branches on \texttt{jl} (less-than) with two
register operands; \texttt{cmp3.s} compares against an immediate:

\begin{lstlisting}
        cmp $3,%r9              # compare %r9 against the immediate 3
        jl less
\end{lstlisting}

\subsection*{A for-loop in assembly}

\texttt{for1.s} sums 1 through 10:

\begin{lstlisting}
_start:
        mov $1,%rcx
        mov $0,%r8

loop:
        cmp $11,%rcx
        jge finish
        add %rcx,%r8
        add $1,%rcx
        jmp loop

finish:
        mov %r8,%rdi
        mov $60,%rax
        syscall
\end{lstlisting}

\texttt{\%rcx} is the conventional loop counter. The structure is
``test before the body'' --- a while-style loop. Exit status is 55.

\subsection*{Globals and the wonder of \texttt{\%rip}-relative addressing}

\texttt{for2.s} introduces global variables in the \texttt{.data}
section and walks them through the same loop:

\begin{lstlisting}
        .section .data
num:    .quad 0
num2:   .quad 0

        .section .text
        .global _start

_start:
        mov $1,%rcx

loop:
        cmp $11,%rcx
        jge finish
        add %rcx,num(%rip)
        addq $1,num2(%rip)
        add $1,%rcx
        jmp loop

finish:
        mov num(%rip),%rdi
        mov $60,%rax
        syscall
\end{lstlisting}

The mysterious \texttt{(\%rip)} bit is RIP-relative addressing.
Here's the conceptual picture: \texttt{\%rip} is the instruction
pointer, the address of the instruction currently executing. The
assembler doesn't know where \texttt{num} will live at runtime
(that depends on the loader and ASLR), but it \emph{does} know how
many bytes \texttt{num} is from this very instruction. So
\texttt{num(\%rip)} is really shorthand for ``take the runtime
\texttt{\%rip}, add this assembler-computed offset, that's where
\texttt{num} lives.''

This is the modern, position-independent way to refer to globals.
We'll keep using it.

\section{Lecture 7 --- Globals, \texttt{lea}, Arrays, \texttt{read}/\texttt{write}, and Functions}

A dense one. We made global variables real, used \texttt{lea} to
take addresses, walked an array four ways, did \texttt{read} and
\texttt{write} both in C and in assembly, and wrote our first
function.

\subsection*{Globals, properly}

\texttt{global1.s} is the cleanest possible single-global program:

\begin{lstlisting}
        .section .data
num:    .quad 10                # 64 bit integer (long)

        .section .text
        .global _start

_start:
        addq $10,num(%rip)
        movq num(%rip), %rdi
        movq $60, %rax
        syscall
\end{lstlisting}

A few things to notice:
\begin{itemize}
  \item Globals declare a \emph{size}, not a type. \texttt{.quad} is
    8 bytes, \texttt{.long} is 4, \texttt{.word} is 2,
    \texttt{.byte} is 1.
  \item \texttt{addq} has a \texttt{q} suffix because we're
    operating on 64-bit (quadword) data. The size suffixes
    (\texttt{b}, \texttt{w}, \texttt{l}, \texttt{q}) become
    important when memory operands don't pin down the size.
  \item I write \texttt{.global \_start} because the linker has to
    see that symbol from outside; I do \emph{not} write
    \texttt{.global num} because nothing outside this file needs
    to refer to it.
\end{itemize}

\subsection*{\texttt{lea}: how pointers actually work}

\texttt{lea} (``load effective address'') is the assembly
equivalent of C's \texttt{\&}. \texttt{global2.s}:

\begin{lstlisting}
_start:
        lea num(%rip),%r9
        # %r9 = &num
        addq $10,(%r9)
        # *(%r9) = *(%r9) + 10
        movq (%r9), %rdi
        movq $60, %rax
        syscall
\end{lstlisting}

The parens around a register --- \texttt{(\%r9)} --- mean
``dereference,'' the same as the C \texttt{*} operator. So
\texttt{addq \$10, (\%r9)} adds 10 to the 8 bytes at the address
held in \texttt{\%r9}. That maps to C as
\texttt{*ptr = *ptr + 10}.

\subsection*{Arrays in assembly}

Three files, one new addressing trick at a time.

\texttt{globalArray1.s} declares an array as a sequence of values
under one label:

\begin{lstlisting}
arr:    .quad 1,2,3,4

_start:
        lea arr(%rip),%r9
        mov (%r9), %rdi         # *arr -> 1
        mov $60, %rax
        syscall
\end{lstlisting}

That gives us element 0. To get other elements, \texttt{globalArray2.s}
uses scale-index addressing:

\begin{lstlisting}
        lea arr(%rip),%r9
        mov $1,%rcx
        mov (%r9,%rcx,8), %rdi
        # *(%r9 + 1*8)
\end{lstlisting}

The form \texttt{(base, index, scale)} computes
\texttt{base + index*scale}, where the scale is 1, 2, 4, or 8.
Since our elements are \texttt{.quad} (8 bytes), the scale is 8.

\texttt{globalArray3.s} sums a 5-element array with this pattern in
a loop:

\begin{lstlisting}
arr:    .quad 1,2,3,4,5

_start:
        lea arr(%rip),%r8
        mov $0,%r9              # accumulator
        mov $0,%rcx             # i = 0

loopStart:
        cmp $5,%rcx
        jge loopEnd
        add (%r8,%rcx,8),%r9
        add $1,%rcx
        jmp loopStart

loopEnd:
        mov %r9,%rdi
        mov $60,%rax
        syscall
\end{lstlisting}

Exit status is 15. This pattern --- base in one register, index in
another, scale equal to element size --- is exactly how C array
indexing compiles.

\subsection*{\texttt{read} and \texttt{write}, first in C}

\texttt{readWrite1.c} uses the raw POSIX \texttt{write} syscall
wrapper:

\begin{lstlisting}[language=C]
#include "stdio.h"
#include "unistd.h"

int main(){
  char hiThere[] = "Hey there, goofuses and goofuxen\n";
  write(1, hiThere, 33);
  return 0;
}
\end{lstlisting}

The arguments are \texttt{(fd, buf, count)}, where \texttt{1} is
stdout. \texttt{readWrite2.c} mirrors it for input:

\begin{lstlisting}[language=C]
char buff[30];
int charsRead = read(0,buff,30);
write(1,buff,charsRead);
\end{lstlisting}

\texttt{0} is stdin. \texttt{read} returns the number of bytes it
actually got, and we feed that into \texttt{write} so we don't
echo back uninitialized garbage at the tail of the buffer.

\subsection*{\texttt{read} and \texttt{write}, in assembly}

\texttt{readWriteASM1.s} writes a fixed string:

\begin{lstlisting}
        .section .data
hiThere: .ascii "Hey there, goofuses and goofuxen\n"

        .section .text
        .global _start

_start:
        mov $1,%rdi
        lea hiThere(%rip),%rsi
        mov $33,%rdx
        mov $1,%rax
        syscall

        mov $60,%rax
        mov $0,%rdi
        syscall
\end{lstlisting}

The Linux x86-64 syscall convention: syscall number in
\texttt{\%rax}, then arguments in \texttt{\%rdi}, \texttt{\%rsi},
\texttt{\%rdx}, \texttt{\%r10}, \texttt{\%r8}, \texttt{\%r9}.
\texttt{write} is syscall 1, \texttt{(fd, buf, count)}.

\texttt{readWriteASM2.s} chains \texttt{read} and \texttt{write}
through a \texttt{.bss} buffer:

\begin{lstlisting}
        .section .bss
buff:   .skip 100

        .section .text
        .global _start

_start:
        mov $0,%rdi
        lea buff(%rip),%rsi
        mov $100,%rdx
        mov $0,%rax
        syscall
        ## %rax now holds the number of bytes read
        mov $1,%rdi
        lea buff(%rip),%rsi
        mov %rax,%rdx
        mov $1,%rax
        syscall

        mov $60,%rax
        mov $0,%rdi
        syscall
\end{lstlisting}

\texttt{.bss} is zero-initialized scratch storage; \texttt{.skip 100}
reserves 100 bytes. The clever line is \texttt{mov \%rax, \%rdx} ---
the syscall returns its byte-count in \texttt{\%rax}, and we funnel
that into \texttt{\%rdx} as the count argument for \texttt{write},
so we only echo back what was actually typed.

\subsection*{Functions and the calling convention}

The Linux x86-64 calling convention for user-space functions:
\begin{itemize}
  \item Arguments in order: \texttt{\%rdi}, \texttt{\%rsi},
    \texttt{\%rdx}, \texttt{\%rcx}, \texttt{\%r8}, \texttt{\%r9}.
  \item Return value in \texttt{\%rax}.
\end{itemize}
(Syscalls are slightly different --- the fourth argument is
\texttt{\%r10}, not \texttt{\%rcx} --- but for plain function calls,
the list above is what you want.)

\texttt{firstFunc.s} is the first ``real'' function we wrote:

\begin{lstlisting}
sillyAdd:
        mov $0,%rax
        add %rdi,%rax
        add %rsi,%rax
        ret

_start:
        mov $10,%rdi
        mov $20,%rsi
        call sillyAdd

        mov %rax,%rdi
        mov $60,%rax
        syscall
\end{lstlisting}

\texttt{call sillyAdd} pushes the address of the next instruction
onto the stack and jumps; \texttt{ret} pops that address and jumps
back. The stack carries the return address invisibly --- that's the
whole magic trick.

\section{Lecture 8 --- The Stack, Local Variables, \texttt{readInt}/\texttt{writeInt}, and Recursion}

Today we made the call stack real, wrote a function with locals by
hand, linked in pre-written \texttt{readInt} and \texttt{writeInt}
helpers, and warned about stack overflow via runaway recursion.

\subsection*{Logistics}

Coding assignment 2 is due in two weeks. The next discussions are
fluffier so people can really sit with assembly. Quizzes are
unlimited attempts --- if you don't like your score, just try
again.

\subsection*{Local variables and the stack}

Every process gets a chunk of address space laid out roughly:
\begin{itemize}
  \item Stack at the top, growing \emph{downward}.
  \item Heap (\texttt{malloc}-land) growing upward from somewhere
    lower.
  \item Data sections (\texttt{.data}, \texttt{.bss},
    \texttt{.rodata}) below that.
  \item Executable instructions at the bottom.
\end{itemize}
Local variables live on the stack. Each function call gets a
\emph{stack frame} of contiguous bytes. Two registers run the show:
\begin{itemize}
  \item \texttt{\%rsp} --- stack pointer (top of stack, lowest
    address).
  \item \texttt{\%rbp} --- base pointer (bottom of the current
    frame).
\end{itemize}

\subsection*{\texttt{push} and \texttt{pop}}

\texttt{pushpop1.s} demonstrates the two primitives:

\begin{lstlisting}
_start:
        push $10
        pop %rdi

        mov $60,%rax
        syscall
\end{lstlisting}

\texttt{push \$10} subtracts 8 from \texttt{\%rsp} and stores 10
there. \texttt{pop \%rdi} loads from \texttt{(\%rsp)} into
\texttt{\%rdi} and adds 8 back. Exit status: 10.

\subsection*{The function prologue and epilogue}

To carve out local space:
\begin{enumerate}
  \item \texttt{push \%rbp} --- save the caller's base pointer.
  \item \texttt{mov \%rsp, \%rbp} --- new frame starts here.
  \item \texttt{sub \$N, \%rsp} --- reserve $N$ bytes (multiple of
    16 for ABI alignment).
\end{enumerate}
Local $i$ then lives at \texttt{-(8*i)(\%rbp)}. To clean up:
\begin{enumerate}
  \item \texttt{mov \%rbp, \%rsp} --- pop all locals at once.
  \item \texttt{pop \%rbp} --- restore caller's base pointer.
  \item \texttt{ret}.
\end{enumerate}

\texttt{stackfunk1.s} puts it together:

\begin{lstlisting}
sillyFunk:
        push %rbp
        mov %rsp,%rbp
        sub $16,%rsp
        movq $1,-8(%rbp)
        movq $2,-16(%rbp)
        mov $0,%rax
        add -8(%rbp),%rax
        add -16(%rbp),%rax

        mov %rbp,%rsp
        pop %rbp
        ret

_start:
        call sillyFunk
        mov %rax,%rdi
        mov $60,%rax
        syscall
\end{lstlisting}

The C analogue is just:

\begin{lstlisting}[language=C]
long sillyfunk(){
  long num1 = 1;
  long num2 = 2;
  return num1 + num2;
}
\end{lstlisting}

Exit status is 3.

\subsection*{Linking external functions: \texttt{readInt} and \texttt{writeInt}}

\texttt{readInt.s} reads a line from stdin and parses it as an
integer:

\begin{lstlisting}
        .section .bss
buffer: .zero 128

        .section .text
        .global readInt

parseInt:
intRead:
        movb (%rdi,%rcx,1), %r8b
        sub $'0', %r8
        imul $10, %rax
        add %r8, %rax
        inc %rcx
        cmp %rsi, %rcx
        jl intRead
        ret

readInt:
        mov $0, %rax
        mov $0, %rdi
        lea buffer(%rip), %rsi
        mov $128, %rdx
        syscall

        mov %rax, %rbx
        mov $0, %rcx
        mov $0, %rax
        lea buffer(%rip), %rdi
        dec %rbx
        mov %rbx, %rsi

        call parseInt
        ret
\end{lstlisting}

Heads up --- this is one of our recurring ``almost-but-not-quite''
bugs. \texttt{movb (\%rdi, \%rcx, 1), \%r8b} loads one byte into
the low byte of \texttt{\%r8}, but \emph{doesn't clear the upper
bits}. If \texttt{\%r8} happened to have garbage in its upper 56
bits, the subsequent \texttt{sub} and \texttt{add} would compound
that garbage. The corrected idiom uses \texttt{movzbl} (move byte
to long, zero-extended) so the upper bits get zeroed for free. The
version in the assembly guide uses \texttt{movzbl} for exactly
this reason. Worth fixing if you're using this as a starting point
for your own code.

\texttt{readInt} also forgets to push/pop \texttt{\%rbx} around its
use --- \texttt{\%rbx} is callee-saved in the System V ABI, so a
caller is entitled to find it intact when we return. That's another
fix-it-please.

\texttt{writeInt.s} is the mirror image: it takes an integer in
\texttt{\%rdi}, allocates a 32-byte buffer on the stack, fills it
\emph{from the end} by repeatedly dividing by 10 and stashing the
ASCII digit, then writes the result. It handles negatives, has a
special case for \texttt{INT\_MIN} (which can't be negated in
two's complement, so we load $2^{63}$ via \texttt{movq} of a 64-bit
immediate), and uses \texttt{leave} as the shorthand for
``\texttt{mov \%rbp, \%rsp; pop \%rbp}.''

(There's also \texttt{writeIntBad.s} sitting in the directory ---
identical to \texttt{writeInt.s} except the \texttt{.global writeInt}
line is commented out. The point: without \texttt{.global}, the
linker can't find the symbol from another file, and you get an
``undefined reference to writeInt'' error.)

\texttt{echoInt.s} is the smallest possible client:

\begin{lstlisting}
        .section .text
        .extern readInt,writeInt
        .global _start

_start:
        call readInt
        mov %rax,%rdi
        call writeInt

        mov $0,%rdi
        mov $60,%rax
        syscall
\end{lstlisting}

\texttt{.extern} declares that \texttt{readInt} and \texttt{writeInt}
live in another file. Build by assembling each separately and
linking together.

\subsection*{Recursion and stack overflow}

Every recursive call allocates a fresh stack frame. The Linux
default stack is around 4MB --- not infinite. \texttt{overflowPractice.c}
shows what happens when you go too deep:

\begin{lstlisting}[language=C]
int counter = 0;

int badRec(int n){
  int stuff = 1;
  int moreStuff = 2;
  counter++;
  if(n == 0){
    return 0;
  }
  else{
    printf("%d\n",counter);
    return badRec(n+1)+stuff+moreStuff;
  }
}

int main(){
  badRec(1000);
  return 0;
}
\end{lstlisting}

Notice the bug: \texttt{badRec(n+1)} instead of \texttt{badRec(n-1)}.
This will never reach the base case --- it just keeps allocating
frames until the stack collides with the heap and the kernel kills
the process. This is the canonical ``stack overflow.''

\section{Lecture 9 --- Processes and the First Look at Threads}

A pivot day: assembly roundup, the discussion problem in detail,
and then the start of concurrency --- processes, fork, and a teaser
on threads.

\subsection*{Discussion 4 walkthrough}

The discussion was: write a C program that reads two numbers and
adds them, then use \texttt{gcc -S} to find the addition in the
assembly and change it to a multiplication. \texttt{discussionDiscussion.c}
is the C version:

\begin{lstlisting}[language=C]
int main(){
  int num1; // -8(%rbp)
  int num2; // -12(%rbp)

  printf("Enter a number: ");
  scanf("%d",&num1);
  printf("Enter another: ");
  scanf("%d",&num2);

  int num3 = num1 + num2; // num3 is -4(%rbp)
  printf("Ta-da! %d\n",num3);
  return 0;
}
\end{lstlisting}

In \texttt{discussionDiscussion.s} the patch is at the addition
site --- the original was an \texttt{addl}, and the discussion
modification turns it into:

\begin{lstlisting}
        movl    -8(%rbp), %edx
        movl    -12(%rbp), %eax
        imul    %edx, %eax
        movl    %eax, -4(%rbp)
\end{lstlisting}

There's also \texttt{dd2.s} where I tried it with the unary form
\texttt{mull \%edx} (which multiplies \texttt{\%eax} by the operand
implicitly --- that's the historical x86 form, where the result
spans \texttt{\%edx:\%eax}). Both work; \texttt{imul} with two
operands is cleaner.

\subsection*{Cross-platform notes on calling convention}

\begin{itemize}
  \item Windows on x86-64 has the \emph{same} instruction set ---
    that's set by the chip --- but a different calling convention:
    arguments go in \texttt{\%rcx}, \texttt{\%rdx}, \texttt{\%r8},
    \texttt{\%r9} for the first four, the rest on the stack, and
    there's a 32-byte ``shadow space'' the caller has to reserve.
  \item Modern Macs are on ARM64 (Apple Silicon). The instruction
    set is totally different, but the calling convention is largely
    Unix-shaped.
\end{itemize}
The compiler is what bridges all this --- one language, different
codegen for each target.

\subsection*{Processes}

A \emph{process} is the OS-level unit of ``a running program.''
The OS bundles together the executable code, register state,
file handles, mapped memory, and so on into one process data
structure.

A ``thread of execution'' is the actual entity scheduled onto a CPU
core. A process has at least one.

\subsection*{\texttt{fork}: the OG concurrency primitive}

\texttt{fork1.c} is the simplest possible:

\begin{lstlisting}[language=C]
#include "stdio.h"
#include "stdlib.h"
#include "unistd.h"

int main(){
  printf("How's it going?\n");

  fork(); // create a new process that's a CLONE of the original

  printf("Hi there\n");
  return 0;
}
\end{lstlisting}

You'll see ``How's it going?'' once and ``Hi there'' twice, because
after \texttt{fork()} there are two processes and both run the rest
of \texttt{main}.

\texttt{fork2.c} distinguishes parent from child by inspecting the
return value of \texttt{fork}:

\begin{lstlisting}[language=C]
int num1 = 10;
pid_t pid = fork();

if(pid != 0){
  num1++;
  printf("I'm the parent. Here's my number  %d\n",num1);
}
else{
  num1--;
  printf("I'm the child. Here's my number %d\n",num1);
}
\end{lstlisting}

\texttt{fork} returns the child's PID in the parent and 0 in the
child. The two processes have \emph{independent} copies of
\texttt{num1} --- forking clones the entire memory space, so the
parent and child are now operating on different bytes that just
happen to have the same name.

\texttt{fork3.c} shows what happens when you fork in a loop ---
you get $2^n$ processes:

\begin{lstlisting}[language=C]
int main(){
  for(int i=0; i< 4;i++){
    fork();
    printf("I'm Mr. %d here at index %d\n",getpid(),i);
  }
  return 0;
}
\end{lstlisting}

After 4 iterations there are 16 processes flooding stdout. This is
how runaway forks become fork bombs in real life.

(One thing missing here that we'll fix later: there's no
\texttt{wait()} in the parent, so children can become zombies until
\texttt{init} reaps them. Add \texttt{wait(NULL);} after \texttt{fork()}
in real code.)

\subsection*{Threads, briefly}

The crucial difference between processes and threads: when you
make a new thread, you do \emph{not} clone everything. A thread
gets its own registers (including its own \texttt{\%rip} and stack)
but \emph{shares the address space} with all the other threads in
the process. That's the source of all the power and all the pain.

\texttt{thread1.c} is the minimum:

\begin{lstlisting}[language=C]
#include <stdio.h>
#include <pthread.h>

void* thread_function(void* arg) {
    (void)arg;
    printf("Thread executing\n");
    return NULL;
}

int main() {
    pthread_t thread;
    pthread_create(&thread, NULL, thread_function, NULL);
    pthread_join(thread, NULL);
    printf("Thread joined\n");
    return 0;
}
\end{lstlisting}

\texttt{pthread\_create} takes a thread handle, an attributes
object (NULL for ``default''), the function to run, and a
\texttt{void*} argument to pass to it. \texttt{pthread\_join} waits
for the thread to finish.

\texttt{thread2.c} is the sneak preview of why concurrency is hard:

\begin{lstlisting}[language=C]
int num = 0;

void* thread_function(void* arg) {
  for(int i=0; i<100; i++){
    int num2 = num;
    num2 = num2+1;
    num = num2;
  }
  return NULL;
}

int main() {
  pthread_t thread[100];
  for(int i=0;i<100;i++){
    pthread_create(&thread[i], NULL, thread_function, NULL);
  }
  for(int i=0; i<100; i++){
    pthread_join(thread[i], NULL);
  }
  printf("And the total is: %d\n",num);
  return 0;
}
\end{lstlisting}

In math-land, that should print 10000 (100 threads each
incrementing 100 times). In real life, you'll often see less,
because the read--modify--write of \texttt{num} isn't atomic and
two threads can step on each other. That's a \emph{race condition},
and we'll spend the next couple of lectures fixing it.

\section{Lecture 11 --- Function Pointers, Threads in Earnest, and Mutexes}

(There was no Lecture 10 directory.) Today: function pointers,
because \texttt{pthread\_create} demands one; \texttt{void*} as the
universal pointer type; threads with arguments and return values;
race conditions; mutexes; and packaging shared state in a struct.

\subsection*{Function pointers}

A function pointer is just an address you can call. The syntax is
ugly but the idea is simple. \texttt{funcpoint1.c} writes a
\texttt{map}:

\begin{lstlisting}[language=C]
int incWrongly(int a){
  return (a + 2);
}

void map(int arr[], int sz, int (*func)(int)){
  for(int i=0;i<sz;i++){
    arr[i] = func(arr[i]);
  }
}
\end{lstlisting}

The parameter \texttt{int (*func)(int)} is read as ``a pointer to
a function that takes an \texttt{int} and returns an \texttt{int}.''
You call it just like a regular function: \texttt{func(arr[i])}.

\texttt{funcpoint2.c} adds a two-argument \texttt{reduce}:

\begin{lstlisting}[language=C]
int reduce(int arr[], int sz,
           int (*func)(int,int),
           int base){
  int acc = base;
  for(int i=0;i<sz;i++){
    acc = func(acc,arr[i]);
  }
  return acc;
}
\end{lstlisting}

Function pointers are how we get higher-order functions in C ---
which we need because \texttt{pthread\_create}'s signature
\emph{requires} a function pointer.

\subsection*{\texttt{void} versus \texttt{void*}}

\texttt{voidArith.c} demonstrates that \texttt{void*} arithmetic
behaves like \texttt{char*} arithmetic --- it advances one byte per
\texttt{+1}:

\begin{lstlisting}[language=C]
int arr[5];
void* weirdo = arr;
printf("weirdo: %p\n",weirdo);
printf("weirdo+1: %p\n",weirdo+1);
\end{lstlisting}

This is a GCC extension; standard C says \texttt{void*} arithmetic
is undefined. But the practical takeaway: \texttt{void} is ``the
type of nothing,'' and \texttt{void*} is ``the universal donor of
pointer types'' --- anything can be cast to or from \texttt{void*},
which is how C does generic programming.

\subsection*{Threads, with arguments and return values}

\texttt{threads1.c} is the bare-bones version we already saw last
time. \texttt{threads2.c} passes an argument:

\begin{lstlisting}[language=C]
void* func(void* arg){
  printf("Hi I'm a thread and I'm holding: %d\n",
         *((int*)arg));
  return NULL;
}

int main(){
  pthread_t thread;
  int arg = 10;

  pthread_create(&thread,NULL,func,&arg);
  pthread_join(thread,NULL);
  return 0;
}
\end{lstlisting}

The cast inside the worker --- \texttt{*((int*)arg)} --- is the
universal pattern: take the \texttt{void*}, lie about its type,
dereference. Don't pass pointers to short-lived stack variables
across thread boundaries; use heap-allocated or longer-lived
storage.

\texttt{threads3.c} returns a value:

\begin{lstlisting}[language=C]
void* func(void* arg){
  int convArg = *((int*)arg);
  int* ret = malloc(sizeof(int));
  *ret = 2*convArg;
  return ret;
}

int main(){
  pthread_t thread;
  int arg = 10;
  int* res;

  pthread_create(&thread,NULL,func,&arg);
  pthread_join(thread,(void**)&res);

  printf("The return value from the thread is: %d\n",*res);
  free(res);
  return 0;
}
\end{lstlisting}

\texttt{pthread\_join} takes a \texttt{void**} (yes, two stars) so
it can write the worker's return value back into a caller-side
pointer. The thread \texttt{malloc}s its return so it survives past
the thread's death; the main thread \texttt{free}s it.

\subsection*{Race conditions: \texttt{threadRace.c}}

\begin{lstlisting}[language=C]
int counter = 0;

void* incer(void* arg){
  int temp = counter;
  sleep(rand()%3+1); // simulated work
  counter = temp+1;
  return NULL;
}

int main(){
  pthread_t threads[100];
  for(int i=0; i<100;i++){
    pthread_create(threads+i,NULL,incer,NULL);
  }
  for(int i=0; i<100;i++){
    pthread_join(threads[i],NULL);
  }
  printf("The counter is: %d\n",counter);
  return 0;
}
\end{lstlisting}

The \texttt{sleep} is fake-work that makes the race \emph{much}
more likely to happen. With 100 threads each doing a non-atomic
read-modify-write of \texttt{counter}, you'll often see a final
value of 1 --- everyone read 0, slept, then wrote 1.

\subsection*{Mutexes: \texttt{threadRace2.c}}

\begin{lstlisting}[language=C]
int counter = 0;
pthread_mutex_t mutex;

void* incer(void* arg){
  pthread_mutex_lock(&mutex);

  int temp = counter;
  sleep(rand()%3+1);
  counter = temp+1;
  pthread_mutex_unlock(&mutex);
  return NULL;
}

int main(){
  pthread_t threads[5];
  pthread_mutex_init(&mutex,NULL);

  for(int i=0; i<5;i++){
    pthread_create(threads+i,NULL,incer,NULL);
  }
  for(int i=0; i<5;i++){
    pthread_join(threads[i],NULL);
  }
  printf("The counter is: %d\n",counter);
  pthread_mutex_destroy(&mutex);
  return 0;
}
\end{lstlisting}

A mutex is a thing where exactly one thread can hold the lock at a
time --- everyone else \texttt{lock}s and waits. Wrap your critical
section between \texttt{lock} and \texttt{unlock}, and the
read-modify-write becomes atomic from the rest of the world's
perspective. Now the count is reliably 5.

\subsection*{Bundling state with its lock: \texttt{threadRace3.c}}

\begin{lstlisting}[language=C]
struct guardedCounter {
  int counter;
  pthread_mutex_t mutex;
};

void* incer(void* arg){
  struct guardedCounter* g = (struct guardedCounter*)arg;

  pthread_mutex_lock(&g->mutex);
  int temp = g->counter;
  sleep(rand()%3+1);
  g->counter = temp+1;
  pthread_mutex_unlock(&g->mutex);
  return NULL;
}
\end{lstlisting}

The argument is a \texttt{void*} that's secretly a
\texttt{struct guardedCounter*}. Putting the mutex inside the
struct means anyone with access to the data also has access to its
lock --- the protection travels with the resource. This is the
shape of nearly all real concurrent code.

For more details and a longer worked discussion of these patterns,
see \texttt{concurrencyGuide.org} in the repo.

\section{File I/O --- Files, File Descriptors, \texttt{argv}, and a Line Editor}

(This was our deep dive into file I/O, which fits naturally between
the threading lectures and the signal lectures since it doesn't
involve concurrency itself.)

\subsection*{Everything is a file}

In Unix, ``file'' is a wildly general abstraction. Regular files,
directories, devices, terminal emulators, named pipes, and sockets
all show up to your program through the same set of operations
(\texttt{open}, \texttt{read}, \texttt{write}, \texttt{close}). The
kernel maintains a per-process table of open file objects, and a
\emph{file descriptor} is just an integer index into that table.

\subsection*{File descriptors and \texttt{open}}

\texttt{fd1.c} is the minimal open/close:

\begin{lstlisting}[language=C]
#include <stdio.h>
#include <unistd.h>
#include <fcntl.h>

int main(){
  char fname[] = "Agenda.org";
  int fd = open(fname,O_RDONLY);

  printf("stdin: %d\n",STDIN_FILENO);
  printf("stdout: %d\n",STDOUT_FILENO);
  printf("stderr: %d\n",STDERR_FILENO);

  if(fd >= 0){
    printf("Hey we opened the file %s and its file descriptor is: %d\n",fname,fd);
  }
  else {
    printf("We couldn't open the file\n");
  }
  close(fd);
  return 0;
}
\end{lstlisting}

\texttt{STDIN\_FILENO}, \texttt{STDOUT\_FILENO}, and
\texttt{STDERR\_FILENO} are the macros for 0, 1, 2 respectively.
The first \texttt{open} call gets you fd 3 (next free slot).

\subsection*{Command-line arguments: \texttt{argv}}

\texttt{fd3.c} is the canonical \texttt{argc}/\texttt{argv} demo:

\begin{lstlisting}[language=C]
int main(int argc, char* argv[]){
  printf("%d arguments were passed in\n",argc);
  for(int i=0; i<argc; i++){
    printf("The %dth argument is: %s\n",i,argv[i]);
  }
  return 0;
}
\end{lstlisting}

\texttt{argv[0]} is conventionally the program name; the actual
arguments start at \texttt{argv[1]}. \texttt{argc} is the count.

\texttt{fd2.c} chains those: it tries to \texttt{open} every argv
entry, prints the resulting fd (or an error), and closes everything
at the end. Notice that \texttt{argv[0]} is the program name itself
and we happily try to \texttt{open} it --- which works! The binary
is just another file.

\subsection*{Reading bytes: \texttt{read}}

\texttt{read1.c} reads up to 1024 bytes from a file:

\begin{lstlisting}[language=C]
char readData[1024];
int bytesRead = read(fds[1],readData,1024);
printf("We read in %d out of the file %s\n", bytesRead, argv[1]);
\end{lstlisting}

\texttt{read2.c} reads twice in a row, demonstrating that the
file's read position advances --- the second \texttt{read} picks
up where the first left off.

\texttt{read3.c} loops until \texttt{read} returns 0 (end of file)
or negative (error):

\begin{lstlisting}[language=C]
char readData[10];
int bytesRead = 0;
while((bytesRead = read(fds[1],readData,10)) > 0){
  printf("We read in %d out of the file %s\n", bytesRead, argv[1]);
}
\end{lstlisting}

This is the canonical ``read until done'' loop.

\subsection*{Writing our own \texttt{cat}}

\texttt{mycat.c} concatenates all the files named on the command
line and writes them to stdout:

\begin{lstlisting}[language=C]
int main(int argc, char* argv[]){
  int numFiles = argc - 1;
  int* fds = malloc(numFiles * sizeof(int));

  for(int i=0; i < numFiles; i++){
    fds[i] = open(argv[i+1],O_RDONLY);
  }

  for(int i=0; i < numFiles; i++){
    char readData[1024];
    int bytesRead = 0;
    while((bytesRead = read(fds[i],readData,1024)) > 0){
      write(STDOUT_FILENO,readData,bytesRead);
    }
  }
  for(int i=0; i< argc; i++){
    close(fds[i]);
  }
  return 0;
}
\end{lstlisting}

That's basically what \texttt{cat} does. (Subtle bug in the close
loop: it goes from 0 to \texttt{argc}, but \texttt{fds} is only
\texttt{numFiles = argc - 1} long. The last \texttt{close(fds[argc-1])}
reads off the end of the array. Not a crash in practice, but worth
fixing.)

\subsection*{The \texttt{FILE*} world: \texttt{fopen}, \texttt{fgets}, \texttt{fputs}}

The \texttt{stdio} ``\texttt{FILE*}'' interface is a buffered
wrapper over file descriptors. \texttt{file1.c}:

\begin{lstlisting}[language=C]
FILE* ourFile = fopen(argv[1],"r+");

char line[1024];
while( fgets(line,1024,ourFile) != NULL){
  printf("%s",line);
}
\end{lstlisting}

\texttt{fgets} reads a line at a time --- much friendlier than
raw \texttt{read} when you actually have line-structured data.

\texttt{weirdcopy.c} uses \texttt{fopen} and \texttt{fputs} to copy
a file line-by-line.

\subsection*{Parsing CSVs: \texttt{petsDB.c} and \texttt{petsDB2.c}}

\texttt{petsDB.c} parses one line of a comma-separated
\texttt{name,species,age} file:

\begin{lstlisting}[language=C]
struct petData {
  char name[50];
  char species[50];
  int age;
};

void readPet(FILE* f, struct petData* p){
  fscanf(f,"%49[^,],%49[^,],%d",p->name,p->species,&p->age);
}
\end{lstlisting}

The \texttt{\%49[\^{},]} format means ``read up to 49 chars that
aren't a comma.'' \texttt{petsDB2.c} extends this into a full
loop, using the fact that \texttt{fscanf} returns the number of
fields it parsed --- when that drops below 3 we know we're done:

\begin{lstlisting}[language=C]
while(readPet(petFile,&p) == 3){
  printf("%s is a %s and is %d years old\n",p.name,p.species,p.age);
}
\end{lstlisting}

\subsection*{Directories with \texttt{opendir}/\texttt{readdir}}

\texttt{directory1.c} lists regular files in a directory:

\begin{lstlisting}[language=C]
DIR* dir = opendir(path);
struct dirent *entry;
while ((entry = readdir(dir)) != NULL) {
  if(entry->d_type == DT_REG){
    printf("%s\n", entry->d_name);
  }
}
closedir(dir);
\end{lstlisting}

\texttt{readdir} returns one entry per call; \texttt{NULL} means
we're done. The \texttt{d\_type} field tells you regular file vs.
directory vs. symlink vs. socket vs. etc.

\texttt{openDir.c} is similar but listed as part of the same
exercise.

\subsection*{File metadata with \texttt{lstat}/\texttt{fstat}}

\texttt{lstatTest.c} (and the identical \texttt{fstatTest.c}) builds
a tiny version of \texttt{ls -l}, decoding the mode bits to print
file type and permissions:

\begin{lstlisting}[language=C]
struct stat file_stat;
lstat(argv[1], &file_stat);

switch (file_stat.st_mode & S_IFMT) {
 case S_IFREG:  printf("regular file\n"); break;
 case S_IFDIR:  printf("directory\n"); break;
 case S_IFLNK:  printf("symbolic link\n"); break;
 // ...
}

printf("%c", file_stat.st_mode & S_IRUSR ? 'r' : '-');
printf("%c", file_stat.st_mode & S_IWUSR ? 'w' : '-');
printf("%c", file_stat.st_mode & S_IXUSR ? 'x' : '-');
// ...

printf("Size of file in bytes: %ld\n", (long) file_stat.st_size);
\end{lstlisting}

\texttt{lstat} doesn't follow symlinks (it tells you about the link
itself), \texttt{fstat} works on an already-open fd, and plain
\texttt{stat} follows symlinks. The mode bits are a bitfield ---
mask with \texttt{S\_IRUSR} etc. to read individual permissions.

\subsection*{Putting it together: a line editor}

\texttt{lineEditor.c} is the grand finale. It loads a file into a
\texttt{char**} (an array of lines), then offers an interactive
menu to edit, delete, or insert lines, and writes the result back
on quit. The interesting bits:

\begin{lstlisting}[language=C]
void cleanUp(FILE* f, char** ls, int nL){
  ftruncate(fileno(f),0);
  fseek(f, 0, SEEK_SET);
  for(int i=0; i < nL; i++){
    fputs(ls[i],f);
    free(ls[i]);
  }
  free(ls);
  fclose(f);
}

void insLine(int line, char** ls, int* nL){
  char* newLine = malloc(linesize*sizeof(char));
  // shift everything from line..end down by one
  for(int i = *nL; i > line; i--){
    ls[i] = ls[i-1];
  }
  ls[line] = newLine;
  (*nL)++;
}
\end{lstlisting}

\texttt{ftruncate} on the underlying fd zeroes the file; then we
\texttt{fseek} to the start and write everything back. The line
shifting is the classic ``insert into an array'' move --- walk
backwards from the end so you don't overwrite data you haven't
copied yet.

There's also a tiny demo \texttt{cursing.c} for \texttt{ncurses}
(initialize, print, refresh, wait for keypress, end) --- ncurses is
the standard library for terminal-based UIs and is what you'd
build a more polished editor on top of.

\section{Lecture 13 --- Signals and Signal Handlers}

Today: exceptions, processes, virtual memory, and signals --- the
mechanism by which the OS interrupts your normal control flow.

\subsection*{Exceptions, broadly}

The OS interrupts your program in three flavors:
\begin{itemize}
  \item \emph{Hardware interrupts}: external signals to the CPU,
    e.g.\ a key press. Handled by a kernel-maintained interrupt
    descriptor table.
  \item \emph{Faults}: things like segfaults (serious) and page
    faults (often routine --- ``go fetch this page from RAM'').
  \item \emph{Traps and aborts}: divide-by-zero, illegal
    instruction, and similar.
\end{itemize}

\subsection*{What's in a process?}

Each process bundles together:
\begin{itemize}
  \item The actual executable code.
  \item Saved register state (what the kernel snapshots on a context
    switch).
  \item Open file descriptors.
  \item Mapped memory pages.
\end{itemize}
The \emph{scheduler} is the part of the OS that decides which
process gets the CPU next. A \emph{context switch} is the (relatively
expensive) operation of saving one process's registers and
restoring another's.

\emph{Virtual memory} gives every process the illusion that it owns
the entire address space. The MMU translates virtual addresses to
physical ones on the fly, which is how processes are isolated from
each other. When working sets exceed RAM, the system starts
\emph{swapping} pages to disk and everything grinds to a halt.

\subsection*{\texttt{getpid} and \texttt{pidLoop.c}}

\texttt{pidLoop.c} is the simplest possible ``find my PID and yell
about it'' loop:

\begin{lstlisting}[language=C]
int main(){
 pid_t myPID = getpid();
 while(1){
  printf("I'm mr. %d\n",myPID);
  sleep(1);
 }
 return 0;
}
\end{lstlisting}

Run it in one terminal, take the PID it prints, then in another
terminal you can \texttt{kill -9 <pid>} or \texttt{kill -USR1 <pid>}
to send signals at it. That's the setup for everything below.

\subsection*{Common signals}

\begin{itemize}
  \item SIGINT --- Ctrl-C
  \item SIGFPE --- divide by zero (and friends)
  \item SIGKILL --- ``you're going to die, no questions, no handler''
  \item SIGSEGV --- segfault
  \item SIGALRM --- alarm timer
  \item SIGSTOP / SIGTSTP --- stop / Ctrl-Z
  \item SIGUSR1, SIGUSR2 --- user-defined, no system meaning
\end{itemize}

\subsection*{Catching SIGINT: \texttt{uninterruptable.c}}

\begin{lstlisting}[language=C]
void uninterruptable(int signum){
  (void)signum;
  write(1,"\nOuch!\n",6);
}

int main(){
  signal(SIGINT, uninterruptable);

  while(1){
    printf("I'm Mr. %d and I cannot be stopped!\n\n",getpid());
    sleep(1);
  }
  return 0;
}
\end{lstlisting}

Install a handler with \texttt{signal(SIGNAL, handler\_function)};
now Ctrl-C prints ``Ouch!'' instead of killing the process. To
actually quit you need a second terminal and \texttt{kill -9}, or
use SIGTERM, etc.

(Note: inside a signal handler, \texttt{printf} is not safe ---
it's not async-signal-safe and can deadlock if the signal arrived
mid-\texttt{printf}. Use \texttt{write} on the raw fd, as we do
here.)

\subsection*{Custom signals: \texttt{sigusr1.c} and \texttt{sigusr2.c}}

\texttt{sigusr1.c} listens for two user-defined signals:

\begin{lstlisting}[language=C]
void sigh_andler(int signum){
  if(signum == SIGUSR1){
    char msg[] = "I guess we're done?\n";
    write(1,msg,strlen(msg));
    exit(0);
  }
  else if(signum == SIGUSR2){
    char msg2[] = "Did you say something?\n";
    write(1,msg2,strlen(msg2));
  }
}

int main(){
  signal(SIGUSR1,sigh_andler);
  signal(SIGUSR2,sigh_andler);

  while(1){
    printf("I'm Mr. %d and I'm just minding my own business...\n",getpid());
    sleep(1);
  }
}
\end{lstlisting}

From another terminal: \texttt{kill -USR1 <pid>} makes it exit;
\texttt{kill -USR2 <pid>} makes it talk back.

\texttt{sigusr2.c} introduces \texttt{volatile sig\_atomic\_t}
state shared between handler and main loop:

\begin{lstlisting}[language=C]
volatile sig_atomic_t state = 0;

void sigh_andler(int signum){
  if(signum == SIGUSR1){
    state +=1;
  }
  else if(signum == SIGUSR2){
    state -= 1;
  }
}

int main(){
  signal(SIGUSR1,sigh_andler);
  signal(SIGUSR2,sigh_andler);

  while(state < 5){
    printf("I'm Mr. %d and I'm just minding my own business...\n",getpid());
    sleep(1);
  }
}
\end{lstlisting}

\texttt{volatile} tells the compiler ``don't cache this in a
register, reload it from memory every time'' --- otherwise an
optimizer might decide \texttt{state < 5} can never change inside
the loop and turn the whole thing into \texttt{while(true)}.
\texttt{sig\_atomic\_t} is the type guaranteed to be readable and
writable atomically with respect to signals.

\subsection*{Signaling between fork-children: \texttt{sigusr3.c}}

\begin{lstlisting}[language=C]
int main(){
  signal(SIGUSR1,sigh_andler);
  signal(SIGUSR2,sigh_andler);

  pid_t pid = fork();
  if(pid == 0){
    // child: ping the parent every two seconds
    for(int i=0; i<5;i++){
      kill(getppid(), SIGUSR1);
      sleep(2);
    }
  }

  while(state < 5){
    printf("I'm Mr. %d and I'm just minding my own business...\n",getpid());
    sleep(1);
  }
}
\end{lstlisting}

The child uses \texttt{getppid()} (parent PID) and
\texttt{kill(pid, SIG)} to send SIGUSR1 to the parent five times.
The parent's handler bumps \texttt{state} each time, and once
\texttt{state >= 5} the parent exits.

\subsection*{An incremental game with signals: \texttt{forkIncremental.c}}

\begin{lstlisting}[language=C]
volatile sig_atomic_t income = 0;
volatile sig_atomic_t gen1 = 0;
volatile sig_atomic_t gen2 = 0;
volatile sig_atomic_t gen3 = 0;

void heartbeat(int signumb){
  printMenu();
  gen2 += gen3;
  gen1 += gen2;
  income += gen1;
}

int main(){
  pid_t pinger = fork();

  if(pinger == 0){
    while(1){
      pid_t parent = getppid();
      sleep(1);
      kill(parent,SIGUSR1);
    }
  }

  signal(SIGUSR1,heartbeat);

  printMenu();
  while(1){
    int choice;
    scanf("%d",&choice);
    switch(choice){
    case 1: income++; break;
    case 2:
      if(income >= 10){ income -= 10; gen1++; }
      break;
    // ...
    }
  }
}
\end{lstlisting}

A child process pings the parent once per second with SIGUSR1; the
parent's handler advances the game state. Meanwhile the parent's
main loop reads keyboard input. This is a real little incremental
``cookie-clicker''-style game using nothing but signals as the
heartbeat. Next lecture we'll redo this same game with threads
instead, which is a much cleaner pattern.

\section{Lecture 14 --- Threads in Earnest, and Reworking the Incremental Game}

The thread review day: re-doing the incremental game with threads
and mutexes instead of signals, plus some standalone thread review
examples and a teaser on \texttt{sigaction}.

\subsection*{Logistics for the assignment}

The next coding assignment is parallel processing: do some
operation in parallel on a big malloc'd array (hundreds of
thousands or millions of elements). Pick a fixed thread count (say
10), and have each thread grab a fixed-size chunk (say 1000
elements) at a time, using a shared counter (protected by a mutex)
to track the next chunk. Don't divide the array by the number of
threads --- the chunk size should be a constant.

\subsection*{Thread review: \texttt{threadreview1.c}}

\begin{lstlisting}[language=C]
void* threadWorker(void* arg){
  char* msg = (char*)arg;
  usleep(rand()%15000);
  printf("%s", msg);
  return NULL;
}

int main(){
  pthread_t thread1, thread2, thread3;

  char msg1[] = "Hi there\n";
  char msg2[] = " how are you?\n";
  char msg3[] = " I hope this message came through\n";

  pthread_create(&thread1,NULL,threadWorker,msg1);
  pthread_create(&thread2,NULL,threadWorker,msg2);
  pthread_create(&thread3,NULL,threadWorker,msg3);

  pthread_join(thread1,NULL);
  pthread_join(thread2,NULL);
  pthread_join(thread3,NULL);
}
\end{lstlisting}

Three threads each print a different message after a random short
sleep. The output order is unpredictable, which is the point ---
threads run when the scheduler decides, not when you wrote the
\texttt{pthread\_create}.

(The \texttt{usleep} takes microseconds, so \texttt{rand()\%15000}
sleeps up to 15 ms.) After joining all three, we then re-create
\texttt{thread1} again to show that you can reuse the
\texttt{pthread\_t} variable once the previous use is joined.

\subsection*{Mutex review: \texttt{threadreview2.c}}

\begin{lstlisting}[language=C]
char* msgs[] = {"Hi there ","how are you? ","I hope this message came through"};
int counter = 0;
pthread_mutex_t mut;

void* threadWorker(void* arg){
  sleep(rand()%3+1);

  int temp;
  pthread_mutex_lock(&mut);
  temp = counter;
  counter++;
  pthread_mutex_unlock(&mut);
  printf("%s", msgs[temp]);
  return NULL;
}
\end{lstlisting}

Each thread grabs the lock, snapshots and increments a shared
counter, releases the lock, then uses the snapshotted value to
pick which message to print. The lock makes the
``read-and-increment'' indivisible.

\subsection*{Little simulation: \texttt{test1.c}}

\texttt{test1.c} is a sketch of an agent simulation:

\begin{lstlisting}[language=C]
enum lilguyState { eating, sleeping, moving, vibing };

struct lilguyData {
  int id;
  enum lilguyState state;
};

sig_atomic_t simulationRunning = 1;

void handler(int signum){
  char msg[] = "Alright everyone clean up and go!\n";
  simulationRunning = 0;
  write(1,msg,strlen(msg));
}

void* threadWorker(void* arg){
  struct lilguyData dat = *((struct lilguyData*) arg);
  while(simulationRunning){
    printf("I'm Miss %d and I'm ",dat.id);
    printState(dat.state);
    dat.state = rand()%4;
    sleep(rand()%3+1);
  }
  return NULL;
}
\end{lstlisting}

Five ``little guy'' threads each randomly cycle through states
until SIGINT (Ctrl-C) flips \texttt{simulationRunning} to 0 and
they all clean up and exit.

\subsection*{Incremental game, take two: \texttt{threadIncremental.c}}

\begin{lstlisting}[language=C]
pthread_mutex_t mut;
int income = 0, gen1 = 0, gen2 = 0, gen3 = 0;

void* heartbeat(void* arg){
  while(1){
    pthread_mutex_lock(&mut);
    printMenu();
    gen2 += gen3;
    gen1 += gen2;
    income += gen1;
    pthread_mutex_unlock(&mut);
    sleep(1);
  }
  return NULL;
}

int main(){
  pthread_t pulse;
  pthread_mutex_init(&mut,NULL);
  pthread_create(&pulse,NULL,heartbeat,NULL);

  while(1){
    int choice;
    scanf("%d",&choice);
    switch(choice){
    case 1:
      pthread_mutex_lock(&mut);
      income++;
      pthread_mutex_unlock(&mut);
      break;
    // ...
    }
  }
}
\end{lstlisting}

Same game as last time, but now the heartbeat is a worker thread
(not a signal-driven side-effect) and every read or write of the
shared state is wrapped in mutex lock/unlock. This is dramatically
cleaner than the SIGUSR1-driven version, and is how you'd actually
build something like this in production.

\subsection*{Adding a clean shutdown: \texttt{threadIncremental2.c}}

\begin{lstlisting}[language=C]
volatile sig_atomic_t toContinue = 1;

void cleanup(int signum){
  toContinue = 0;
}

// ...
struct sigaction sigsettings;
sigsettings.sa_handler = cleanup;
sigsettings.sa_flags = 0; // makes scanf not block
sigaction(SIGINT,&sigsettings,NULL);

while(toContinue){
  int choice;
  scanf("%d",&choice);
  // ...
}
\end{lstlisting}

We add a SIGINT handler that flips \texttt{toContinue} to 0, and
both the heartbeat loop and main loop check that flag. The catch:
plain \texttt{signal()} sets \texttt{SA\_RESTART} by default on
Linux, which means \texttt{scanf} will silently retry after the
signal --- you'd have to type something to break out. Using
\texttt{sigaction} with \texttt{sa\_flags = 0} disables that, so
\texttt{scanf} returns with an error on signal and the loop checks
\texttt{toContinue}.

\texttt{threadIncremental3.c} and
\texttt{threadIncremental3\_sigaction.c} are tiny variations on
the same theme: one with the plain \texttt{signal()} call (which
will get stuck in \texttt{scanf}), one with \texttt{sigaction}
(which won't), and a version where the cleanup handler also writes
a final ``make one last purchase'' message. The diff between the
two versions is the whole pedagogical point: \texttt{signal()} is
the old, slightly-broken interface; \texttt{sigaction()} is what
you actually want for real code.

\section{Sockets --- Networks, the TCP/IP Stack, and Three Echo Servers}

This was our networking day. Final project handling has moved to
GitHub at \texttt{https://github.com/clarissalittler/cs201}; the
\texttt{sockets-lecture/} directory has the working examples.

\subsection*{The TCP/IP model in 5 minutes}

\begin{itemize}
  \item \emph{Application} layer: HTTP, DNS, SSH, BitTorrent ---
    things built on top of the protocols.
  \item \emph{Transport} layer: TCP (reliable, ordered) and UDP
    (best-effort, unordered, fast). TCP guarantees every byte
    arrives and arrives in order; UDP just lobs packets and hopes.
  \item \emph{Internet} layer: IP addresses (IPv4, 32-bit, four
    dotted octets; IPv6, 128-bit) and routing.
  \item \emph{Link} layer: the physical connection --- ethernet,
    wi-fi.
\end{itemize}

\subsection*{Sockets and network byte order}

A \emph{socket} is the OS abstraction for a network endpoint.
Surprise: it's just a file descriptor, so once you've opened one
you can \texttt{read}/\texttt{write} on it like any other file.

Network byte order is \emph{big-endian} (most-significant byte
first), which is the opposite of x86-64. So before you stuff a port
number or IP address into the socket structures, you call
\texttt{htons} (host-to-network short) or \texttt{htonl}
(host-to-network long) to byte-swap.

\subsection*{Echo server, simplest form}

\texttt{echo\_server\_simplistic.c} is the bare minimum:

\begin{lstlisting}[language=C]
int server_fd = socket(AF_INET, SOCK_STREAM, 0);

int optval = 1;
setsockopt(server_fd, SOL_SOCKET, SO_REUSEADDR,
           &optval, sizeof(optval));

struct sockaddr_in server_addr;
memset(&server_addr, 0, sizeof(server_addr));
server_addr.sin_family = AF_INET;
server_addr.sin_addr.s_addr = htonl(INADDR_ANY);
server_addr.sin_port = htons(port);

bind(server_fd, (struct sockaddr *)&server_addr, sizeof(server_addr));
listen(server_fd, 10);

while (1) {
    struct sockaddr_in client_addr;
    socklen_t client_len = sizeof(client_addr);
    int client_fd = accept(server_fd,
                           (struct sockaddr *)&client_addr,
                           &client_len);

    char buffer[1024];
    ssize_t bytes_received;
    while ((bytes_received = recv(client_fd, buffer,
                                  sizeof(buffer), 0)) > 0) {
        send(client_fd, buffer, bytes_received, 0);
    }
    close(client_fd);
}
\end{lstlisting}

The five-step recipe:
\begin{enumerate}
  \item \texttt{socket(AF\_INET, SOCK\_STREAM, 0)} --- TCP socket
    over IPv4.
  \item Fill out a \texttt{sockaddr\_in} (family, IP, port) and
    \texttt{bind} it to the socket.
  \item \texttt{listen} marks the socket as accepting connections.
  \item \texttt{accept} blocks until a connection comes in,
    returning a new fd for that client.
  \item \texttt{recv} and \texttt{send} are the socket flavors of
    \texttt{read} and \texttt{write}.
\end{enumerate}

The \texttt{SO\_REUSEADDR} option lets you restart the server
without waiting for the kernel's TIME\_WAIT timeout, which is a
huge quality-of-life thing during development.

You can connect to it with \texttt{ncat localhost 8080} (or any
port you bound) and watch your typing get echoed back.

\subsection*{Echo server with error handling: \texttt{echo\_server.c}}

\texttt{echo\_server.c} is the same logic but with proper
\texttt{perror} calls on every step that can fail, and
\texttt{inet\_ntop} to log the client's IP and port:

\begin{lstlisting}[language=C]
char client_ip[INET_ADDRSTRLEN];
inet_ntop(AF_INET, &client_addr.sin_addr, client_ip, sizeof(client_ip));
printf("Connection from %s:%d\n", client_ip, ntohs(client_addr.sin_port));
\end{lstlisting}

\subsection*{Multi-client echo server: \texttt{echo\_server\_threaded.c}}

The simplistic version handles exactly one client at a time. To
handle many in parallel, we spin up a thread per connection:

\begin{lstlisting}[language=C]
void *handle_client(void *arg) {
    int client_fd = *(int *)arg;
    free(arg);

    char buffer[1024];
    ssize_t bytes_received;
    while ((bytes_received = recv(client_fd, buffer,
                                  sizeof(buffer), 0)) > 0) {
        send(client_fd, buffer, bytes_received, 0);
    }
    close(client_fd);
    return NULL;
}

int main(int argc, char *argv[]) {
    // ... socket/bind/listen as before ...
    while (1) {
        int client_fd = accept(server_fd,
                               (struct sockaddr *)&client_addr,
                               &client_len);

        int *client_fd_ptr = malloc(sizeof(int));
        *client_fd_ptr = client_fd;

        pthread_t thread;
        pthread_create(&thread, NULL, handle_client, client_fd_ptr);
        pthread_detach(thread);
    }
}
\end{lstlisting}

Two things to note. First, we \texttt{malloc} a single \texttt{int}
to pass the fd into the thread --- if we passed \texttt{\&client\_fd}
directly, the next iteration would clobber it before the thread
read it. Second, we \texttt{pthread\_detach} so we don't have to
\texttt{join} --- the thread cleans up after itself when it
returns. (We don't have a meaningful return value to wait on.)

\subsection*{The grand finale: \texttt{chat\_server.c}}

A multi-client \emph{broadcast} server: every connected client sees
every other client's messages. The data structure:

\begin{lstlisting}[language=C]
typedef struct {
    int fd;
    int active;
    char ip[INET_ADDRSTRLEN];
    int port;
} Client;

Client clients[MAX_CLIENTS];
pthread_mutex_t clients_mutex = PTHREAD_MUTEX_INITIALIZER;
\end{lstlisting}

A fixed-size array of clients, protected by a single global mutex.
The broadcast routine grabs the lock and writes to every active
client except the sender:

\begin{lstlisting}[language=C]
void broadcast(char *message, int sender_fd) {
    pthread_mutex_lock(&clients_mutex);
    for (int i = 0; i < MAX_CLIENTS; i++) {
        if (clients[i].active && clients[i].fd != sender_fd) {
            send(clients[i].fd, message, strlen(message), 0);
        }
    }
    pthread_mutex_unlock(&clients_mutex);
}
\end{lstlisting}

This is everything we've built all term, all in one place: file
descriptors, threads, mutexes, structured shared state. Try
connecting two terminals to the same server with \texttt{ncat} and
watch the messages cross-pollinate.

For the deeper details (especially around partial reads, message
framing, and what \texttt{recv} \emph{actually} returns at TCP
boundaries), see \texttt{socketsGuide.org} in the repo.

\section{Appendix: Pointer Pyrotechnics (\texttt{pointerExtras/})}

These were extra examples I dropped into the repo for practice ---
not tied to a specific lecture, but worth a read.

\subsection*{\texttt{pointymcpointface.c}: pointers as a substitute for pass-by-reference}

\begin{lstlisting}[language=C]
void worldsWorstInc(int* p){
  *p = *p + 1;
}

int main(){
  int num = 0;
  printf("%p\n",(void*)&num);
  int* np = &num;

  printf("The variable with address %p has a value of %d\n", (void*)np, *np);

  worldsWorstInc(np);
  worldsWorstInc(np);

  printf("The variable with address %p now has a value of %d\n", (void*)np, *np);
}
\end{lstlisting}

The two reasons pointers exist in C: dynamic memory
(\texttt{malloc}/\texttt{free}), and faking pass-by-reference. This
is the latter, distilled.

\subsection*{\texttt{multiarray.c}: dynamic 2D arrays as pointers-to-pointers}

\begin{lstlisting}[language=C]
int** arr = malloc(10*sizeof(int*));
for(int i = 0; i<10; i++){
  arr[i] = malloc(10*sizeof(int));
}

for(int i=0; i < 10; i++){
  for(int j=0; j < 10; j++){
    arr[i][j] = 100*(i+1) + j;
  }
}
\end{lstlisting}

A 10x10 matrix as an array of 10 pointers, each pointing to a row
of 10 ints. The double-subscript \texttt{arr[i][j]} expands to
\texttt{*(*(arr+i)+j)} --- two dereferences and two pointer
arithmetics. Note this isn't contiguous in memory like a true 2D
array (\texttt{int arr[10][10]}); the rows can be allocated wherever.

\subsection*{\texttt{pointstack.c}: a linked-list stack with pointer-to-pointer for the head}

\begin{lstlisting}[language=C]
struct Node {
  int v;
  struct Node* next;
};

void pushStack(struct Node** s, int v){
  struct Node* h = malloc(sizeof(struct Node));
  h->v = v;
  h->next = *s;
  *s = h;
}
\end{lstlisting}

The function takes \texttt{struct Node**} (not \texttt{struct
Node*}) because it needs to mutate the caller's head pointer.
This is the same ``pass an address so the function can modify the
variable'' move from \texttt{worldsWorstInc}, just one level of
indirection deeper. If you internalize \texttt{**} here, every
later use of pointer-to-pointer (e.g.\ \texttt{argv}, the
\texttt{void**} in \texttt{pthread\_join}) becomes much less
mysterious.

\end{document}
